\section{本讲主线：正确性靠编程模型，性能靠硬件模型}
Lecture 6 接在 GPU 课之后，目标是把“知道 GPU 很快”推进到“知道怎么测、怎么定位、怎么写 kernel”。核心循环是：benchmark 看到端到端时间，profile 看到时间花在哪里，然后用更贴近硬件的数据组织和 kernel 写法修正瓶颈。
\begin{importantbox}{本讲的闭环}
编程模型给正确性，硬件模型给性能，benchmark/profiling 给证据，Triton 给可控的 block-level 实现。没有测量就改 kernel，基本是在猜。
\end{importantbox}

\begin{knowledgebox}{课堂提示：正确性与性能需要两套模型}
老师在总结中把职责分得很清楚：PyTorch、Triton 与 PTX 等 programming model 让程序语义可表达、可验证；SM、warp、occupancy、bank conflict 与 memory hierarchy 决定实现能否高效。Benchmark 用于观察 scaling，profiler 用于看真实执行了什么、花了多久；Triton 则把思考单位提升到 thread block，形成“从 HBM 读入、片上处理或 fusion、再写回”的统一路径。
\end{knowledgebox}

\begin{knowledgebox}{术语消化：GPU kernel 基础词汇}
\begin{itemize}
    \item \textbf{Kernel}：在 GPU 上执行的函数。
    \item \textbf{HBM}：High Bandwidth Memory，GPU 的大容量全局显存，带宽高但仍远慢于片上 register/shared memory。
    \item \textbf{Thread block / CTA}：一组线程，被调度到同一个 SM，共享 shared memory。
    \item \textbf{Warp}：通常 32 个 lockstep 执行的 threads；分支不同会产生 control divergence。
    \item \textbf{Occupancy}：SM 上活跃 warps/blocks 的比例，受 registers、shared memory、block size 限制。
    \item \textbf{Bank conflict}：多个 threads 同时访问 shared memory 同一个 bank，访问被串行化。
    \item \textbf{Memory coalescing}：warp 中线程访问连续 HBM cache line，合并成高效 memory transaction。
    \item \textbf{Tiling}：把大矩阵分块搬到 shared memory 中复用。
    \item \textbf{Fusion}：把多个操作合并进一个 kernel，减少 HBM 往返。
\end{itemize}
\end{knowledgebox}

\begin{knowledgebox}{符号说明：kernel 性能公式}
本讲出现的 \(N\) 表示向量或矩阵维度，\(M,N,K\) 在 matmul 中分别表示 \(A\in\mathbb{R}^{M\times K}\)、\(B\in\mathbb{R}^{K\times N}\)、\(C\in\mathbb{R}^{M\times N}\) 的维度；\(\mathrm{FLOPs}\) 是浮点操作数，bytes moved 是 HBM 读写量。Arithmetic intensity \(I=\mathrm{FLOPs}/\mathrm{bytes}\) 越高，越容易让 compute units 忙起来。
\end{knowledgebox}
